%!TEX root = forallxyyc.tex
\part{Dedução natural para a LPO}
\label{ch.NDFOL}
\addtocontents{toc}{\protect\mbox{}\protect\hrulefill\par}

\chapter{Regras básicas para a LPO}\label{s:BasicFOL}

A linguagem da LPO utiliza todos os conectivos da LVF. Assim, as provas na LPO usarão todas as regras básicas e derivadas de \cref{ch.NDTFL}. Também usaremos as noções teóricas da prova (em particular, o símbolo \textquotedblleft $\proves$\textquotedblright{}) introduzidas ali. No entanto, também precisaremos de algumas novas regras básicas para reger os quantificadores e o sinal de identidade.


\section{Eliminação universal}\label{s:forallElim}

Da afirmação de que tudo é~$F$, você pode inferir que qualquer coisa particular é~$F$. Dê-lhe um nome; ela é~$F$. Portanto, o seguinte deve ser permitido:
\begin{fitchproof}
	\hypo{a}{\forall x\,\atom{R}{x,x,d}}\PR
	\have{c}{\atom{R}{a,a,d}} \Ae{a}
\end{fitchproof}
Obtivemos a linha~$2$ eliminando o quantificador universal e substituindo cada ocorrência de \textquotedblleft $x$\textquotedblright{} por \textquotedblleft $a$\textquotedblright{}. Do mesmo modo, o seguinte deve ser permitido:
\begin{fitchproof}
	\hypo{a}{\forall x\,\atom{R}{x,x,d}}\PR
	\have{c}{\atom{R}{d,d,d}} \Ae{a}
\end{fitchproof}
Aqui obtivemos a linha~$2$ eliminando o quantificador universal e substituindo cada ocorrência de \textquotedblleft $x$\textquotedblright{} por~\textquotedblleft $d$\textquotedblright{}. Poderíamos ter feito o mesmo com qualquer outro nome que quiséssemos.

Isso motiva a regra de eliminação universal ($\forall$E):

\factoidbox{
\begin{fitchproof}
	\have[m]{a}{\forall \metav{x}\,\metav{A}(\ldots \metav{x} \ldots \metav{x}\ldots)}
	\have[\ ]{c}{\metav{A}(\ldots \metav{c} \ldots \metav{c}\ldots)} \Ae{a}
\end{fitchproof}}

A notação aqui foi introduzida em \cref{s:TruthFOL}. O ponto é
que você pode obter qualquer \emph{instância de substituição} de uma fórmula
quantificada universalmente: substitua cada ocorrência livre da variável
quantificada por qualquer nome que desejar.

Devemos enfatizar que (como ocorre com toda regra de eliminação) você só pode aplicar a regra $\forall$E quando o quantificador universal for o operador lógico principal. Portanto, o seguinte é \emph{proibido}:
\begin{fitchproof}
	\hypo{a}{\forall x\,\atom{B}{x} \eif \atom{B}{k}}\PR
	\have{c}{\atom{B}{b} \eif \atom{B}{k}}\by{tentativa indevida de invocar $\forall$E}{a}
\end{fitchproof}
Isso é ilegítimo, pois \textquotedblleft $\forall x$\textquotedblright{} não é o operador lógico principal na linha~$1$. (Se precisar lembrar por que esse tipo de inferência deve ser proibido, releia \cref{s:MoreMonadic}.)

\section{Introdução existencial}\label{s:existsIntro}

Da afirmação de que alguma coisa particular é~$F$, você pode inferir que algo é~$F$. Portanto, devemos permitir:
\begin{fitchproof}
	\hypo{a}{\atom{R}{a,a,d}}\PR
	\have{b}{\exists x\, \atom{R}{a,a,x}} \Ei{a}
\end{fitchproof}
Aqui substituímos o nome \textquotedblleft $d$\textquotedblright{} por uma variável \textquotedblleft $x$\textquotedblright{} e então quantificamos existencialmente essa variável. Do mesmo modo, poderíamos ter permitido:
\begin{fitchproof}
	\hypo{a}{\atom{R}{a,a,d}}\PR
	\have{c}{\exists x\, \atom{R}{x,x,d}} \Ei{a}
\end{fitchproof}
Aqui substituímos ambas as ocorrências do nome \textquotedblleft $a$\textquotedblright{} por uma variável e então generalizamos existencialmente. Mas não precisamos substituir \emph{ambas} as ocorrências de um nome por uma variável: se Narciso ama a si mesmo, então há alguém que ama Narciso. Portanto, também permitimos:
\begin{fitchproof}
	\hypo{a}{\atom{R}{a,a,d}}\PR
	\have{d}{\exists x\, \atom{R}{x,a,d}} \Ei{a}
\end{fitchproof}
Aqui substituímos \emph{uma} ocorrência do nome \textquotedblleft $a$\textquotedblright{} por uma variável e então generalizamos existencialmente. Essas observações motivam nossa regra de introdução, embora, para explicá-la, precisemos introduzir alguma notação nova.

Quando $\metav{A}$ é uma sentença que contém o nome $\metav{c}$,
podemos enfatizar isso escrevendo
\textquotedblleft $\metav{A}(\ldots \metav{c} \ldots \metav{c}\ldots)$\textquotedblright{}. Escreveremos
\textquotedblleft $\metav{A}(\ldots \metav{x} \ldots \metav{c}\ldots)$\textquotedblright{} para indicar qualquer fórmula obtida substituindo
\emph{algumas ou todas} as ocorrências do nome $\metav{c}$ pela
variável~$\metav{x}$ (desde que $\metav{c}$, onde quer que seja
substituído, não ocorra no escopo de um quantificador que
ligue~$\metav{x}$). Com isso, nossa regra de introdução
$\exists$I é:
\factoidbox{
\begin{fitchproof}
	\have[m]{a}{\metav{A}(\ldots \metav{c} \ldots \metav{c}\ldots)}
	\have[\ ]{c}{\exists \metav{x}\,\metav{A}(\ldots \metav{x} \ldots \metav{c}\ldots)} \Ei{a}
\end{fitchproof}
%\metav{x} must not occur in $\metav{A}(\ldots \metav{c} \ldots
%\metav{c}\ldots)$
}
% The constraint is included to guarantee that any application of the rule yields a  sentence of FOL. Thus the following is allowed:
% \begin{fitchproof}
% 	\hypo{a}{\atom{R}{a,a,d}}
% 	\have{d}{\exists x\, \atom{R}{x,a,d}} \Ei{a}
% 	\have{e}{\exists y \exists x\, \atom{R}{x,y,d}} \Ei{d}
% \end{fitchproof}
% But this is banned:
% \begin{fitchproof}
% 	\hypo{a}{\atom{R}{a,a,d}}
% 	\have{d}{\exists x\, \atom{R}{x,a,d}} \Ei{a}
% 	\have{e}{\exists x\, \exists x\, \atom{R}{x,x,d}}\by{naughty attempt to invoke $\exists$I}{d}
% \end{fitchproof}
% since the expression on line~3 contains clashing variables, and so is not a sentence of FOL.

\section{Domínios vazios}
A prova a seguir combina nossas duas novas regras para quantificadores:
	\begin{fitchproof}
		\hypo{a}{\forall x\, \atom{F}{x}}\PR
		\have{in}{\atom{F}{a}}\Ae{a}
		\have{e}{\exists x\, \atom{F}{x}}\Ei{in}
	\end{fitchproof}
Esta poderia ser uma prova ruim? Se algo existe, então certamente podemos inferir que algo é~$F$ a partir do fato de que tudo é~$F$. Mas e se \emph{nada} existir? Então é certamente verdade, de modo vacuamente, que tudo é~$F$; no entanto, disso não se segue que algo seja~$F$, pois não há nada que possa \emph{ser}~$F$. Portanto, se afirmamos que \textquotedblleft $\exists x\,\atom{F}{x}$\textquotedblright{} decorre de \textquotedblleft $\forall x\,\atom{F}{x}$\textquotedblright{} apenas por lógica, então estamos afirmando que, apenas por lógica, há algo em vez de nada. Isso pode parecer um pouco estranho.

Na verdade, já estamos comprometidos com essa estranheza. Em \cref{s:FOLBuildingBlocks}, estipulamos que os domínios na LPO devem ter pelo menos um membro. Em seguida, definimos uma validade (da LPO) como uma sentença verdadeira em toda interpretação. Como \textquotedblleft $\exists x\,x=x$\textquotedblright{} será verdadeira em toda interpretação, isso também teve o efeito de estipular que é uma questão de lógica haver algo em vez de nada.

Como está longe de ser claro que a lógica deva nos dizer que deve haver algo em vez de nada, podemos estar trapaceando um pouco aqui.

Se nos recusarmos a trapacear, porém, pagaremos um preço alto. Eis três coisas às quais queremos nos apegar:
	\begin{compactlist}
\item $\forall x\,\atom{F}{x} \proves \atom{F}{a}$: afinal, essa era a regra $\forall$E.
\item $\atom{F}{a} \proves \exists x\,\atom{F}{x}$: afinal, essa era a regra $\exists$I.
\item a capacidade de copiar e colar provas: afinal, o raciocínio funciona juntando muitos passos pequenos em cadeias bastante longas.
	\end{compactlist}
Se conseguirmos o que queremos nos três aspectos, teremos de aceitar que $\forall x\,\atom{F}{x} \proves \exists x\,\atom{F}{x}$. Assim, se conseguirmos o que queremos nos três aspectos, o próprio sistema de provas nos dirá que há algo em vez de nada. E, se nos recusarmos a aceitar isso, teremos de abandonar uma das três coisas às quais queremos nos apegar!

Antes de começar a pensar em qual abandonar, talvez devêssemos perguntar \emph{quanto} estamos trapaceando. É verdade que isso pode dificultar debates teológicos sobre por que há algo em vez de nada. Mas, no restante do tempo, estaremos perfeitamente bem. Talvez devêssemos simplesmente considerar que nosso sistema de provas (e a LPO, de modo mais geral) tem um alcance ligeiramente limitado. Se algum dia quisermos permitir a possibilidade de \emph{nada}, teremos de procurar um sistema de provas mais complicado. Mas, enquanto estivermos contentes em ignorar essa possibilidade, nosso sistema de provas estará perfeitamente em ordem. (Assim como a estipulação de que todo domínio deve conter pelo menos um objeto.)

\section{Introdução universal}\label{s:forallIntro}

Suponha que você tenha mostrado, para cada coisa particular, que ela é $F$ (e que não há outras coisas a considerar). Então estaria justificado afirmar que tudo é $F$. Isso motiva a seguinte regra de prova. Se você estabeleceu cada uma das instâncias de substituição de \textquotedblleft $\forall x\,\atom{F}{x}$\textquotedblright{}, pode inferir \textquotedblleft $\forall x\,\atom{F}{x}$\textquotedblright{}.

Infelizmente, essa regra seria completamente inutilizável. Estabelecer cada instância de substituição exigiria provar \textquotedblleft $\atom{F}{a}$\textquotedblright{}, \textquotedblleft $\atom{F}{b}$\textquotedblright{}, \dots, \textquotedblleft $\atom{F}{j_2}$\textquotedblright{}, \dots, \textquotedblleft $\atom{F}{r_{79002}}$\textquotedblright{}, \ldots, e assim por diante. De fato, como há infinitos nomes na LPO, esse processo nunca chegaria ao fim. Portanto, jamais poderíamos aplicar essa regra. Precisamos ser um pouco mais engenhosos ao formular nossa regra para introduzir a quantificação universal.

Uma solução pode ser inspirada pela consideração de:
$$\forall x\,\atom{F}{x} \therefore \forall y\,\atom{F}{y}$$
Esse argumento deveria ser \emph{obviamente} válido. Afinal, a variação alfabética deveria ser uma questão de gosto, sem consequência lógica. Mas como nosso sistema de provas poderia refletir isso? Suponha que comecemos uma prova assim:
\begin{fitchproof}
	\hypo{x}{\forall x\, \atom{F}{x}} \PR
	\have{a}{\atom{F}{a}} \Ae{x}
\end{fitchproof}
Provamos \textquotedblleft $\atom{F}{a}$\textquotedblright{}. E, naturalmente, nada nos impede de usar a mesma justificativa para provar \textquotedblleft $\atom{F}{b}$\textquotedblright{}, \textquotedblleft $\atom{F}{c}$\textquotedblright{}, \ldots, \textquotedblleft $\atom{F}{j_2}$\textquotedblright{}, \ldots, \textquotedblleft $\atom{F}{r_{79002}}$\textquotedblright{}, \dots, e assim por diante até ficarmos sem espaço, tempo ou paciência. Refletindo sobre isso, porém, vemos que há uma maneira de provar $\atom{F}{\metav{c}}$, para qualquer nome~$\metav{c}$. E, se podemos fazê-lo para \emph{qualquer} coisa, certamente devemos poder dizer que \textquotedblleft $F$\textquotedblright{} é verdadeiro de \emph{tudo}. Isso nos justifica inferir \textquotedblleft $\forall y\,\atom{F}{y}$\textquotedblright{}, assim:
\begin{fitchproof}
	\hypo{x}{\forall x\, \atom{F}{x}}\PR
	\have{a}{\atom{F}{a}} \Ae{x}
	\have{y}{\forall y\, \atom{F}{y}} \Ai{a}
\end{fitchproof}
A ideia crucial aqui é que \textquotedblleft $a$\textquotedblright{} era apenas um nome \emph{arbitrário}. Não havia nada de especial nele — poderíamos ter escolhido qualquer outro nome — e ainda assim a prova estaria correta. Essa ideia crucial motiva a regra de introdução universal ($\forall$I):
\factoidbox{
\begin{fitchproof}
	\have[m]{a}{\metav{A}(\ldots \metav{c} \ldots \metav{c}\ldots)}
	\have[\ ]{c}{\forall \metav{x}\,\metav{A}(\ldots \metav{x} \ldots \metav{x}\ldots)} \Ai{a}
\end{fitchproof}
em que o nome $\metav{c}$ não pode ocorrer em uma premissa ou em uma
hipótese não descarregada.%\\
%\metav{x} must not occur in $\metav{A}(\ldots \metav{c} \ldots
%\metav{c}\ldots)$
}

Em \cref{s:MainLogicalOperatorQuantifier}, introduzimos uma convenção
para indicar \emph{instâncias de substituição}: dissemos que, se
$\metav{A}(\ldots \metav{x} \ldots \metav{x}\ldots)$ é uma fórmula,
então $\metav{A}(\ldots \metav{c} \ldots \metav{c}\ldots)$ é o
resultado de substituir \emph{toda} ocorrência livre de $\metav{x}$ em
$\metav{A}$ por~$\metav{c}$. Para formular nossa regra $\forall$I, usamos
a mesma convenção, mas na direção oposta: se
$\metav{A}(\ldots \metav{c} \ldots \metav{c}\ldots)$ é uma fórmula,
então $\metav{A}(\ldots \metav{x} \ldots \metav{x}\ldots)$ é o
resultado de substituir \emph{toda} ocorrência de~$\metav{c}$ nela
por~$\metav{x}$. (Isso só é permitido enquanto $\metav{c}$ não
ocorrer no escopo de um quantificador que ligue~$\metav{x}$ em~$\metav{A}$.)

Há duas coisas importantes às quais devemos prestar atenção ao aplicar a
regra $\forall$I. Se não formos cuidadosos ao escolher o
nome~$\metav{c}$ e substituí-lo por~$\metav{x}$, corremos o risco de
fazer algumas inferências incorretas.

A restrição da regra $\forall$I exige que $\metav{c}$ não
ocorra em uma premissa ou em uma hipótese não descarregada. Essa restrição
garante que estejamos sempre raciocinando em um nível suficientemente geral.
%\footnote{Recall from \cref{s:BasicTFL} that we are treating `$\ered$' as a canonical contradiction. But if it were the canonical contradiction as involving some \emph{constant}, it might interfere with the constraint mentioned here. To avoid such problems, we will treat `$\ered$' as a canonical contradiction \emph{that involves no particular names}.}
Para ver a restrição em ação, considere este argumento terrível:
	\begin{earg}
\item Todos amam Kylie Minogue.
\item[\texttherefore] Todos amam a si mesmos.
	\end{earg}
Podemos simbolizar esse padrão de inferência obviamente inválido como:
$$\forall x\,\atom{L}{x,k} \therefore \forall x\,\atom{L}{x,x}$$
Agora, suponha que tentássemos oferecer uma prova que justificasse esse argumento:
\begin{fitchproof}
	\hypo{x}{\forall x\, \atom{L}{x,k}}\PR
	\have{a}{\atom{L}{k,k}} \Ae{x}
	\have{y}{\forall x\, \atom{L}{x,x}} \by{tentativa indevida de invocar $\forall$I}{a}
\end{fitchproof}
Isso não é permitido, pois \textquotedblleft $k$\textquotedblright{} já ocorria em uma premissa, a saber, na linha~$1$. O ponto crucial é que, se fizemos alguma suposição sobre o objeto com o qual estamos trabalhando, então não estamos raciocinando de modo suficientemente geral para autorizar $\forall$I.

Embora o nome não possa ocorrer em nenhuma hipótese \emph{não descarregada}, ele pode ocorrer em uma hipótese \emph{descarregada}. Isto é, pode ocorrer em uma subprova que já fechamos. Por exemplo, isto é perfeitamente permitido:
\begin{fitchproof}
	\open
		\hypo{f1}{\atom{G}{d}}\AS
		\have{f2}{\atom{G}{d}}\by{R}{f1}
	\close
	\have{ff}{\atom{G}{d} \eif \atom{G}{d}}\ci{f1-f2}
	\have{zz}{\forall z(\atom{G}{z} \eif \atom{G}{z})}\Ai{ff}
\end{fitchproof}
Isso nos diz que \textquotedblleft $\forall z(\atom{G}{z} \eif \atom{G}{z})$\textquotedblright{} é um \emph{teorema}. E é assim que deve ser.

Se substituirmos apenas \emph{alguns} nomes, e não outros, acabaremos
\textquotedblleft provando\textquotedblright{} coisas absurdas. Por exemplo, considere o argumento:
	\begin{earg}
\item Todos têm a mesma idade que si mesmos.
\item[\texttherefore] Todos têm a mesma idade que Judi Dench.
	\end{earg}
Podemos simbolizar isso da seguinte maneira:
$$\forall x\,\atom{O}{x,x} \therefore \forall x\,\atom{O}{x,d}$$
Mas suponha agora que tentássemos justificar esse argumento terrível com o seguinte:
\begin{fitchproof}
	\hypo{x}{\forall x\, \atom{O}{x,x}}\PR
	\have{a}{\atom{O}{d,d}}\Ae{x}
	\have{y}{\forall x\, \atom{O}{x,d}}\by{tentativa indevida de invocar $\forall$I}{a}
\end{fitchproof}
Felizmente, nossas regras não nos permitem fazer isso: a prova
tentada é proibida, pois não substitui \emph{toda} ocorrência de \textquotedblleft $d$\textquotedblright{}
na linha~$2$ por um~\textquotedblleft $x$\textquotedblright{}.

\section{Eliminação existencial}\label{s:existsElim}

Suponha que saibamos que \emph{algo} é~$F$. O problema é que simplesmente saber isso não nos diz qual coisa é~$F$. Portanto, pareceria que, de \textquotedblleft $\exists x\,\atom{F}{x}$\textquotedblright{}, não podemos concluir imediatamente \textquotedblleft $\atom{F}{a}$\textquotedblright{}, \textquotedblleft $\atom{F}{e_{23}}$\textquotedblright{} ou qualquer outra instância de substituição da sentença. O que podemos fazer?

Suponha que saibamos que algo é~$F$ e que tudo o que é~$F$ também é~$G$. Em um inglês quase natural, poderíamos raciocinar assim:
	\begin{quote}
Como algo é $F$, há alguma coisa particular que é~$F$. Não sabemos nada sobre ela além de que é~$F$, mas, por conveniência, vamos chamá-la de \textquotedblleft Becky\textquotedblright{}. Portanto: Becky é $F$. Como tudo o que é $F$ é~$G$, segue-se que Becky é~$G$. Mas, como Becky é~$G$, segue-se que algo é~$G$. Nada dependeu exatamente de qual objeto Becky era. Portanto, algo é~$G$.
	\end{quote}
Poderíamos tentar capturar esse padrão de raciocínio em uma prova da seguinte maneira:
\begin{fitchproof}
	\hypo{es}{\exists x\, \atom{F}{x}}\PR
	\hypo{ast}{\forall x(\atom{F}{x} \eif \atom{G}{x})}\PR
	\open
		\hypo{s}{\atom{F}{b}}\AS
		\have{st}{\atom{F}{b} \eif \atom{G}{b}}\Ae{ast}
		\have{t}{\atom{G}{b}} \ce{st, s}
		\have{et1}{\exists x\, \atom{G}{x}}\Ei{t}
	\close
	\have{et2}{\exists x\, \atom{G}{x}}\Ee{es,s-et1}
\end{fitchproof}\noindent
Analisando a prova: começamos escrevendo nossas hipóteses. Na linha~$3$, fizemos uma hipótese adicional: \textquotedblleft $\atom{F}{b}$\textquotedblright{}. Essa era apenas uma instância de substituição de \textquotedblleft $\exists x\,\atom{F}{x}$\textquotedblright{}. Com base nessa hipótese, estabelecemos \textquotedblleft $\exists x\,\atom{G}{x}$\textquotedblright{}. Observe que não fizemos nenhuma hipótese \emph{especial} sobre o objeto nomeado por \textquotedblleft $b$\textquotedblright{}; apenas supusemos que ele satisfaz \textquotedblleft $\atom{F}{x}$\textquotedblright{}. Portanto, nada depende de qual objeto ele é. E a linha~$1$ nos disse que \emph{algo} satisfaz \textquotedblleft $\atom{F}{x}$\textquotedblright{}, de modo que nosso padrão de raciocínio foi perfeitamente geral. Podemos descarregar a hipótese específica \textquotedblleft $\atom{F}{b}$\textquotedblright{} e simplesmente inferir \textquotedblleft $\exists x\,\atom{G}{x}$\textquotedblright{} por si só.

Juntando tudo isso, obtemos a regra de eliminação existencial ($\exists$E):
\factoidbox{
\begin{fitchproof}
	\have[m]{a}{\exists \metav{x}\,\metav{A}(\ldots \metav{x} \ldots \metav{x}\ldots)}
	\open
		\hypo[i]{b}{\metav{A}(\ldots \metav{c} \ldots \metav{c}\ldots)}\AS
		\have[j]{c}{\metav{B}}
	\close
	\have[\ ]{d}{\metav{B}} \Ee{a,b-c}
\end{fitchproof}
\metav{c} não pode ocorrer em nenhuma hipótese não descarregada antes da linha $i$\\
\metav{c} não pode ocorrer em $\exists \metav{x}\,\metav{A}(\ldots \metav{x} \ldots \metav{x}\ldots)$\\
\metav{c} não pode ocorrer em \metav{B}}
Assim como na introdução universal, as restrições são extremamente importantes. Para ver por quê, considere o seguinte argumento terrível:
	\begin{earg}
\item Tim Button é professor.
\item Alguém não é professor.
\item[\texttherefore] Tim Button é simultaneamente professor e não professor.
	\end{earg}
Podemos simbolizar este padrão de inferência obviamente inválido da seguinte maneira:
$$\atom{L}{b}, \exists x\, \enot\atom{L}{x} \therefore \atom{L}{b} \eand \enot \atom{L}{b}$$
Agora suponha que tentássemos oferecer uma prova que justificasse esse argumento:
\begin{fitchproof}
	\hypo{f}{\atom{L}{b}}\PR
	\hypo{nf}{\exists x\, \enot \atom{L}{x}}\PR
	\open
		\hypo{na}{\enot \atom{L}{b}}\AS
		\have{con}{\atom{L}{b} \eand \enot \atom{L}{b}}\ai{f, na}
	\close
\ifHTMLtarget
\have{econ1}{\atom{L}{b} \eand \enot \atom{L}{b}}\by{tentativa indevida de invocar $\exists$E }{nf, na-con}

\else
	\have{econ1}{\atom{L}{b} \eand \enot \atom{L}{b}}\by{tentativa indevida}{}
	\have[\ ]{x}{}\by{para invocar $\exists$E }{nf, na-con}
\fi
\end{fitchproof}
A última linha da prova não é permitida. O nome que usamos em nossa instância de substituição de \textquotedblleft $\exists x\, \enot \atom{L}{x}$\textquotedblright{} na linha~$3$, a saber, \textquotedblleft $b$\textquotedblright{}, ocorre na linha~$4$. Isto não seria melhor:
\begin{fitchproof}
	\hypo{f}{\atom{L}{b}}\PR
	\hypo{nf}{\exists x\, \enot \atom{L}{x}}\PR
	\open
		\hypo{na}{\enot \atom{L}{b}}\AS
		\have{con}{\atom{L}{b} \eand \enot \atom{L}{b}}\ai{f, na}
		\have{con1}{\exists x (\atom{L}{x} \eand \enot \atom{L}{x})}\Ei{con}
	\close
\ifHTMLtarget
\have{econ1}{\exists x (\atom{L}{x} \eand \enot \atom{L}{x})}\by{tentativa indevida de invocar $\exists$E }{nf, na-con1}

\else
\have{econ1}{\exists x (\atom{L}{x} \eand \enot \atom{L}{x})}\by{tentativa indevida}{}
	\have[\ ]{x}{}\by{para invocar $\exists$E }{nf, na-con1}
\fi
\end{fitchproof}
A última linha ainda não é permitida. O nome que usamos em nossa instância de substituição de \textquotedblleft $\exists x\, \enot \atom{L}{x}$\textquotedblright{}, a saber, \textquotedblleft $b$\textquotedblright{}, ocorre em uma hipótese não descarregada, isto é, na linha~$1$.

Por fim, considere esta \textquotedblleft prova\textquotedblright{}:
\begin{fitchproof}
	\hypo{nf}{\exists x\, \atom{L}{a, x}}\PR
	\open
		\hypo{na}{\atom{L}{a, a}}\AS
		\have{con1}{\exists x\,\atom{L}{x,x}}\Ei{na}
	\close
\ifHTMLtarget
	\have{econ1}{\exists x\,\atom{L}{x,x}}\by{tentativa indevida de invocar $\exists$E }{nf, na-con1}
\else
	\have{econ1}{\exists x\,\atom{L}{x,x}}\by{tentativa indevida}{}
	\have[\ ]{x}{}\by{para invocar $\exists$E }{nf, na-con1}
\fi
\end{fitchproof}
Aqui o nome \textquotedblleft $a$\textquotedblright{} viola a segunda condição da regra $\exists$E: ele ocorre em \textquotedblleft $\exists x\, \atom{L}{a,x}$\textquotedblright{} na linha~$1$. E precisamos impedir que isto seja uma prova, pois ela é inválida: \textquotedblleft alguém ama a si mesmo\textquotedblright{} não decorre de \textquotedblleft Alex ama alguém\textquotedblright{}.




A lição desta história é a seguinte. \emph{Se você quiser extrair informação de um quantificador existencial, escolha um nome novo para sua instância de substituição.} Assim, poderá garantir que cumpre todas as restrições da regra $\exists$E.

\section{Regras dos quantificadores e quantificação vacuosa}

Há uma questão sutil em nossas regras para quantificadores relacionada a um
aspecto da maneira como especificamos as regras de nossa sintaxe, a saber,
a possibilidade de \emph{quantificação vacuosa} (ver o final de
\cref{s:FreeVars}). Em \cref{s:forallElim,s:existsElim}, usamos
a notação de substituição de \cref{s:TruthFOL}: $\metav{A}(\ldots
\metav{c} \ldots \metav{c} \ldots)$ é a fórmula obtida
substituindo toda ocorrência \emph{livre} de $\metav{x}$ em $\metav{A}$
por~$\metav{c}$. Aqui enfatizamos o termo \textquotedblleft livre\textquotedblright{} por uma razão. Isso significa
que nem sempre você pode substituir \emph{toda} ocorrência de~$\metav{x}$.
Por exemplo, o seguinte estaria incorreto:
\begin{fitchproof}
	\hypo{a}{\forall x(\atom{F}{x} \land \exists x\,\atom{G}{x})}\PR
	\ifHTMLtarget
	\have{e}{\atom{F}{a} \land \exists x\,\atom{G}{a}}\by{tentativa indevida de invocar $\forall$E}{a}
\else
	\have{e}{\atom{F}{a} \land \exists x\,\atom{G}{a}}\by{tentativa indevida}{}
	\have[\ ]{x}{}\by{para invocar $\forall$E}{a}
\fi
\end{fitchproof}
Na linha $2$, formamos incorretamente a instância de substituição:
substituímos o \textquotedblleft $x$\textquotedblright{} tanto em \textquotedblleft $\atom{F}{x}$\textquotedblright{} quanto em \textquotedblleft $\atom{G}{x}$\textquotedblright{}
por~\textquotedblleft $a$\textquotedblright{}, mas somente a primeira ocorrência é livre — a outra
é ligada por~\textquotedblleft $\exists x$\textquotedblright{}. Isso não apenas violaria uma restrição
à substituição, como também permitiria que \textquotedblleft provássemos\textquotedblright{} um argumento inválido.
Como \textquotedblleft $\exists x\,\atom{G}{a}$\textquotedblright{} é um caso de quantificação vacuosa,
é equivalente simplesmente a \textquotedblleft $\atom{G}{a}$\textquotedblright{}. Além disso, a linha~$1$ é
equivalente a \textquotedblleft $\forall x(\atom{F}{x} \land \exists y\,\atom{G}{y})$\textquotedblright{}.
Deve ser fácil ver que isso não implica \textquotedblleft $\atom{F}{a} \land
\atom{G}{a}$\textquotedblright{}. É improvável que você encontre um caso assim (a menos que
tenha um professor ardiloso): seus exercícios provavelmente envolverão apenas
sentenças como \textquotedblleft $\forall x(\atom{F}{x} \land \exists
y\,\atom{G}{y})$\textquotedblright{}, com as variáveis cuidadosamente separadas.


Em \cref{s:existsIntro}, dissemos que usamos \textquotedblleft $\metav{A}(\ldots
\metav{x} \ldots \metav{c}\ldots)$\textquotedblright{} para indicar qualquer fórmula obtida
substituindo \emph{algumas ou todas} as ocorrências do nome $\metav{c}$
pela variável~$\metav{x}$ \emph{desde que~$\metav{c}$,
onde quer que seja substituído, não ocorra no escopo de um quantificador
que ligue~$\metav{x}$}. A restrição enfatizada, que proíbe
a substituição de $\metav{c}$ em alguns casos, é um tanto sutil. Você
provavelmente jamais se encontrará em uma situação em que corra o
risco de violá-la. Mas ela é necessária para tornar a regra correta.
Como permitimos quantificação vacuosa, por exemplo, \textquotedblleft $\exists
x\forall x\,\atom{L}{x,x}$\textquotedblright{} é uma sentença legal da LPO. Ocorre apenas que o primeiro
\textquotedblleft $\exists x$\textquotedblright{} não faz nada, e a sentença
é equivalente simplesmente a \textquotedblleft $\forall x\,\atom{L}{x,x}$\textquotedblright{}. Agora, \textquotedblleft $\forall
x\,\atom{L}{x,x}$\textquotedblright{} \emph{poderia} ser obtida de \textquotedblleft $\forall
x\,\atom{L}{a,x}$\textquotedblright{} substituindo o nome~\textquotedblleft $a$\textquotedblright{} pela variável~\textquotedblleft $x$\textquotedblright{}
se não tivéssemos cuidado e ignorássemos que, neste caso, \textquotedblleft $a$\textquotedblright{} \emph{ocorre}
no escopo de \textquotedblleft $\forall x$\textquotedblright{}. Sem a restrição, o seguinte seria permitido:

\begin{fitchproof}
		\hypo{a}{\forall x\, \atom{L}{a,x}}\PR
		\ifHTMLtarget
		\have{e}{\exists x\, \forall x\, \atom{L}{x,x}}\by{tentativa indevida de invocar $\exists$I}{a}
	\else
		\have{e}{\exists x\, \forall x\, \atom{L}{x,x}}\by{tentativa indevida}{}
		\have[\ ]{x}{}\by{para invocar $\exists$I}{a}
	\fi\end{fitchproof}
Mas, como \textquotedblleft $\exists x\, \forall x\, \atom{L}{x,x}$\textquotedblright{} é equivalente
simplesmente a \textquotedblleft $\forall x\, \atom{L}{x,x}$\textquotedblright{}, isso seria uma inferência inválida.
Para ver por quê, leia \textquotedblleft $L$\textquotedblright{} como \textquotedblleft ama\textquotedblright{}: de
\textquotedblleft Alex ama todos\textquotedblright{} não se segue que todos amem todos.

Da mesma forma, em \cref{s:forallIntro}, dissemos que $\metav{A}(\ldots
\metav{x} \ldots \metav{x}\ldots)$ é o resultado de substituir
\emph{toda} ocorrência de~$\metav{c}$ em $\metav{A}(\ldots \metav{c}
\ldots \metav{c}\ldots)$ por~$\metav{x}$, mas que \emph{isso só é
permitido enquanto $\metav{c}$ não ocorrer no escopo de um
quantificador que ligue~$\metav{x}$ em~$\metav{A}$}. Assim como a restrição
incluída na formulação da regra $\exists$I, isso não costuma acontecer
na prática. Mas eis por que ela é importante. Novamente, como permitimos
quantificação vacuosa, \textquotedblleft $\forall
x\exists x\,\atom{L}{x,x}$\textquotedblright{} é uma sentença legal da LPO. Contudo,
\textquotedblleft $\forall x\exists x\,\atom{L}{x,x}$\textquotedblright{} pode ser obtida de \textquotedblleft $\forall
x\,\atom{L}{a,x}$\textquotedblright{} substituindo o nome~\textquotedblleft $a$\textquotedblright{} pela variável~\textquotedblleft $x$\textquotedblright{}
em toda parte. Isso viola, porém, a restrição que incluímos: \textquotedblleft $a$\textquotedblright{}
ocorre no escopo de \textquotedblleft $\forall x$\textquotedblright{} neste caso. Sem a restrição, o seguinte seria permitido:

\begin{fitchproof}
	\hypo{p}{\forall y\exists x\,\atom{L}{y,x}}\PR
	\have{a}{\exists x\, \atom{L}{a,x}}\Ae{p}
\ifHTMLtarget
	\have{e}{\forall x\, \exists x\, \atom{L}{x,x}}\by{tentativa indevida de invocar $\forall$I}{a}
\else
\have{e}{\forall x\, \exists x\, \atom{L}{x,x}}\by{tentativa indevida}{}
	\have[\ ]{x}{}\by{para invocar $\forall$I}{a}
\fi
\end{fitchproof}
Mas, como \textquotedblleft $\forall x\, \exists x\, \atom{L}{x,x}$\textquotedblright{} é equivalente
simplesmente a \textquotedblleft $\exists x\, \atom{L}{x,x}$\textquotedblright{}, isso seria uma inferência
inválida: \textquotedblleft alguém ama a si mesmo\textquotedblright{} não decorre de
\textquotedblleft todos amam alguém\textquotedblright{}.

\practiceproblems
\problempart
Explique por que estas duas \textquotedblleft provas\textquotedblright{} estão \emph{incorretas}. Além disso, forneça interpretações que invalidem as formas de argumento falaciosas que essas \textquotedblleft provas\textquotedblright{} consagram:
\begin{multicols}{2}
	\begin{fitchproof}
		\hypo{Rxx}{\forall x\, \atom{R}{x,x}}\PR
		\have{Raa}{\atom{R}{a,a}}\Ae{Rxx}
		\have{Ray}{\forall y\, \atom{R}{a,y}}\Ai{Raa}
		\have{Rxy}{\forall x\, \forall y\, \atom{R}{x,y}}\Ai{Ray}
	\end{fitchproof}
	\begin{fitchproof}
		\hypo{AE}{\forall x\, \exists y\, \atom{R}{x,y}}\PR
		\have{E}{\exists y\, \atom{R}{a,y}}\Ae{AE}
		\open
			\hypo{ass}{\atom{R}{a,a}}\AS
			\have{Ex}{\exists x\, \atom{R}{x,x}}\Ei{ass}
		\close
		\have{con}{\exists x\, \atom{R}{x,x}}\Ee{E, ass-Ex}
	\end{fitchproof}
\end{multicols}

\problempart
\label{pr.justifyFOLproof}
As três provas a seguir não têm suas citações (regra e números de linha). Acrescente-as para transformá-las em provas legítimas.
\begin{compactlist}
\item \begin{fitchproof}
\hypo{p1}{\forall x\exists y(\atom{R}{x,y} \eor \atom{R}{y,x})}
\hypo{p2}{\forall x\,\enot \atom{R}{m,x}}
\have{3}{\exists y(\atom{R}{m,y} \eor \atom{R}{y,m})}{}
	\open
		\hypo{a1}{\atom{R}{m,a} \eor \atom{R}{a,m}}
		\have{a2}{\enot \atom{R}{m,a}}{}
		\have{a3}{\atom{R}{a,m}}{}
		\have{a4}{\exists x\, \atom{R}{x,m}}{}
	\close
\have{n}{\exists x\, \atom{R}{x,m}} {}
\end{fitchproof}

\item \begin{fitchproof}
\hypo{1}{\forall x(\exists y\,\atom{L}{x,y} \eif \forall z\,\atom{L}{z,x})}
\hypo{2}{\atom{L}{a,b}}
\have{3}{\exists y\,\atom{L}{a,y} \eif \forall z\atom{L}{z,a}}{}
\have{4}{\exists y\, \atom{L}{a,y}} {}
\have{5}{\forall z\, \atom{L}{z,a}} {}
\have{6}{\atom{L}{c,a}}{}
\have{7}{\exists y\,\atom{L}{c,y} \eif \forall z\,\atom{L}{z,c}}{}
\have{8}{\exists y\, \atom{L}{c,y}}{}
\have{9}{\forall z\, \atom{L}{z,c}}{}
\have{10}{\atom{L}{c,c}}{}
\have{11}{\forall x\, \atom{L}{x,x}}{}
\end{fitchproof}

\item \begin{fitchproof}
\hypo{a}{\forall x(\atom{J}{x} \eif \atom{K}{x})}
\hypo{b}{\exists x\,\forall y\, \atom{L}{x,y}}
\hypo{c}{\forall x\, \atom{J}{x}}
\open
	\hypo{2}{\forall y\, \atom{L}{a,y}}
	\have{3}{\atom{L}{a,a}}{}
	\have{d}{\atom{J}{a}}{}
	\have{e}{\atom{J}{a} \eif \atom{K}{a}}{}
	\have{f}{\atom{K}{a}}{}
	\have{4}{\atom{K}{a} \eand \atom{L}{a,a}}{}
	\have{5}{\exists x(\atom{K}{x} \eand \atom{L}{x,x})}{}
\close
\have{j}{\exists x(\atom{K}{x} \eand \atom{L}{x,x})}{}
\end{fitchproof}
\end{compactlist}

\problempart
\label{pr.BarbaraEtc.proof1}
No \hyperref[pr.BarbaraEtc]{exercício \ref*{s:MoreMonadic}A}, consideramos quinze figuras silogísticas da lógica aristotélica. Forneça
provas para cada uma das formas de argumento. OBS.: será
\emph{muito} mais fácil se você simbolizar (por exemplo)
\textquotedblleft nenhum $F$ é $G$\textquotedblright{} como
\textquotedblleft $\forall x (\atom{F}{x} \eif \enot \atom{G}{x})$\textquotedblright{}.

\

\problempart
\label{pr.BarbaraEtc.proof2}
Aristóteles e seus sucessores identificaram outras formas silogísticas que dependiam da \textquotedblleft importação existencial\textquotedblright{}. Simbolize cada uma dessas formas de argumento na LPO e ofereça provas.
\begin{compactlist}
\item \textbf{Barbari.} Algo é $H$. Todo $G$ é $F$. Todo $H$ é $G$. Logo: algum $H$ é $F$.
\item \textbf{Celaront.} Algo é $H$. Nenhum $G$ é $F$. Todo $H$ é $G$. Logo: algum $H$ não é $F$.
\item \textbf{Cesaro.} Algo é $H$. Nenhum $F$ é $G$. Todo $H$ é $G$. Logo: algum $H$ não é $F$.
\item \textbf{Camestros.} Algo é $H$. Todo $F$ é $G$. Nenhum $H$ é $G$. Logo: algum $H$ não é $F$.
\item \textbf{Felapton.} Algo é $G$. Nenhum $G$ é $F$. Todo $G$ é $H$. Logo: algum $H$ não é $F$.
\item \textbf{Darapti.} Algo é $G$. Todo $G$ é $F$. Todo $G$ é $H$. Logo: algum $H$ é $F$.
\item \textbf{Calemos.} Algo é $H$. Todo $F$ é $G$. Nenhum $G$ é $H$. Logo: algum $H$ não é $F$.
\item \textbf{Fesapo.} Algo é $G$. Nenhum $F$ é $G$. Todo $G$ é $H$. Logo: algum $H$ não é $F$.
\item \textbf{Bamalip.} Algo é $F$. Todo $F$ é $G$. Todo $G$ é $H$. Logo: alguns $H$ são $F$.
\end{compactlist}

\problempart
\label{pr.someFOLproofs}
Para cada uma das afirmações a seguir, forneça uma prova na LPO que mostre
que ela é verdadeira.
\begin{compactlist}
\item $\proves \forall x\,\atom{F}{x} \eif \forall y(\atom{F}{y} \eand \atom{F}{y})$
\item $\forall x(\atom{A}{x}\eif \atom{B}{x}), \exists x\,\atom{A}{x} \proves \exists x\,\atom{B}{x}$
\item $\forall x(\atom{M}{x} \eiff \atom{N}{x}), \atom{M}{a} \eand \exists x\,\atom{R}{x,a} \proves \exists x\,\atom{N}{x}$
\item $\forall x\, \forall y\,\atom{G}{x,y}\proves\exists x\,\atom{G}{x,x}$
\item $\proves\forall x\,\atom{R}{x,x} \eif \exists x\, \exists y\,\atom{R}{x,y}$
\item $\proves\forall y\, \exists x (\atom{Q}{y} \eif \atom{Q}{x})$
\item $\atom{N}{a} \eif \forall x(\atom{M}{x} \eiff \atom{M}{a}), \atom{M}{a}, \enot\atom{M}{b}\proves \enot \atom{N}{a}$
\item $\forall x\, \forall y (\atom{G}{x,y} \eif \atom{G}{y,x}) \proves \forall x\forall y (\atom{G}{x,y} \eiff \atom{G}{y,x})$
\item $\forall x(\enot\atom{M}{x} \eor \atom{L}{j,x}), \forall x(\atom{B}{x}\eif \atom{L}{j,x}), \forall x(\atom{M}{x}\eor \atom{B}{x})\proves \forall x\,\atom{L}{j,x}$
\end{compactlist}

\solutions
\problempart
\label{pr.likes}
Escreva uma chave de simbolização para o argumento a seguir, simbolize-o e prove-o:
\begin{earg}
\item Há alguém que ama todos que amam todos que ela
ama.
\item[\texttherefore] Há alguém que ama a si mesma.
\end{earg}


\problempart
Mostre que cada par de sentenças é interderivável.
\begin{compactlist}
\item $\forall x (\atom{A}{x}\eif \enot \atom{B}{x})$, $\enot\exists x(\atom{A}{x} \eand \atom{B}{x})$
\item $\forall x (\enot\atom{A}{x}\eif \atom{B}{d})$, $\forall x\,\atom{A}{x} \eor \atom{B}{d}$
\item $\exists x\,\atom{P}{x} \eif \atom{Q}{c}$, $\forall x (\atom{P}{x} \eif \atom{Q}{c})$
\end{compactlist}

\solutions
\problempart
\label{pr.FOLequivornot}
Para cada um dos pares de sentenças a seguir: se forem interderiváveis, dê provas que mostrem isso. Se não forem, construa uma interpretação que mostre que não são logicamente equivalentes.
\begin{compactlist}
\item $\forall x\,\atom{P}{x} \eif \atom{Q}{c}, \forall x (\atom{P}{x} \eif \atom{Q}{c})$
\item $\forall x\,\forall y\, \forall z\,\atom{B}{x,y,z}, \forall x\,\atom{B}{x,x,x}$
\item $\forall x\,\forall y\,\atom{D}{x,y}, \forall y\,\forall x\,\atom{D}{x,y}$
\item $\exists x\,\forall y\,\atom{D}{x,y}, \forall y\,\exists x\,\atom{D}{x,y}$
\item $\forall x (\atom{R}{c,a} \eiff \atom{R}{x,a}), \atom{R}{c,a} \eiff \forall x\,\atom{R}{x,a}$
\end{compactlist}

\solutions
\problempart
\label{pr.FOLvalidornot}
Para cada um dos argumentos a seguir: se for válido na LPO, dê uma prova. Se for inválido, construa uma interpretação que mostre que é inválido.
\begin{compactlist}
\item $\exists y\,\forall x\,\atom{R}{x,y} \therefore \forall x\,\exists y\,\atom{R}{x,y}$
\item $\forall x\,\exists y\,\atom{R}{x,y} \therefore  \exists y\,\forall x\,\atom{R}{x,y}$
\item $\exists x(\atom{P}{x} \eand \enot \atom{Q}{x}) \therefore \forall x(\atom{P}{x} \eif \enot \atom{Q}{x})$
\item $\forall x(\atom{S}{x} \eif \atom{T}{a}), \atom{S}{d} \therefore \atom{T}{a}$
\item $\forall x(\atom{A}{x}\eif \atom{B}{x}), \forall x(\atom{B}{x} \eif \atom{C}{x}) \therefore \forall x(\atom{A}{x} \eif \atom{C}{x})$
\item $\exists x(\atom{D}{x} \eor \atom{E}{x}), \forall x(\atom{D}{x} \eif \atom{F}{x}) \therefore \exists x(\atom{D}{x} \eand \atom{F}{x})$
\item $\forall x\,\forall y(\atom{R}{x,y} \eor \atom{R}{y,x}) \therefore \atom{R}{j,j}$
\item $\exists x\,\exists y(\atom{R}{x,y} \eor \atom{R}{y,x}) \therefore \atom{R}{j,j}$
\item $\forall x\,\atom{P}{x} \eif \forall x\,\atom{Q}{x}, \exists x\, \enot\atom{P}{x} \therefore \exists x\, \enot \atom{Q}{x}$
\item $\exists x\,\atom{M}{x} \eif \exists x\,\atom{N}{x}$, $\enot \exists x\,\atom{N}{x}\therefore  \forall x\, \enot \atom{M}{x}$
\end{compactlist}

\chapter{Provas com quantificadores}

Em \cref{s:stratTFL}, discutimos estratégias para construir provas
usando as regras básicas da dedução natural para a LVF. Os mesmos
princípios se aplicam às regras para os quantificadores. Se quisermos provar
uma sentença quantificada $\forall \metav{x}\,
\atom{\metav{A}}{\metav{x}}$ ou $\exists \metav{x}\,
\atom{\metav{A}}{\metav{x}}$, podemos trabalhar para trás justificando a
sentença desejada por $\forall$I ou $\exists$I e tentando encontrar uma
prova da premissa correspondente dessa regra. Para trabalhar para a frente
a partir de uma sentença quantificada, aplicamos $\forall$E ou $\exists$E, conforme
o caso.

Especificamente, suponha que você queira provar $\forall \metav{x}\, \atom{\metav{A}}{\metav{x}}$. Para isso, usando $\forall$I, precisaríamos de uma prova de $\atom{\metav{A}}{\metav{c}}$ para algum nome~$\metav{c}$ que não ocorra em nenhuma hipótese não descarregada. Aplicar a estratégia correspondente, isto é, construir uma prova de $\forall \metav{x}\, \atom{\metav{A}}{\metav{x}}$ trabalhando para trás, consiste em escrever $\atom{\metav{A}}{\metav{c}}$ acima dela e então continuar tentando encontrar uma prova dessa sentença.
\begin{fitchproof}
	\ellipsesline
	\have[n]{n}{\atom{\metav{A}}{\metav{c}}}
	\have{m}{\forall \metav{x}\, \atom{\metav{A}}{\metav{x}}}\Ai{n}
\end{fitchproof}
$\atom{\metav{A}}{\metav{c}}$ é obtida de $\atom{\metav{A}}{\metav{x}}$ substituindo cada ocorrência livre de $\metav{x}$ em $\atom{\metav{A}}{\metav{x}}$ por~$\metav{c}$. Para que isso funcione, $\metav{c}$ deve satisfazer a condição especial. Podemos garantir isso escolhendo sempre um nome que ainda não ocorra na prova construída. (Naturalmente, ele ocorrerá na prova que acabarmos construindo — apenas não em uma hipótese não descarregada na linha~$n+1$.)

Para trabalhar para trás a partir de uma sentença $\exists \metav{x}\, \atom{\metav{A}}{\metav{x}}$, escrevemos de modo semelhante uma sentença acima dela que possa servir de justificativa para uma aplicação da regra $\exists$I, isto é, uma sentença da forma $\atom{\metav{A}}{\metav{c}}$.
\begin{fitchproof}
	\ellipsesline
	\have[n]{n}{\atom{\metav{A}}{\metav{c}}}
	\have{m}{\exists \metav{x}\, \atom{\metav{A}}{\metav{x}}}\Ei{n}
\end{fitchproof}
Isso se parece exatamente com o que faríamos trabalhando para trás a partir de uma sentença quantificada universalmente. A diferença é que, enquanto para $\forall$I temos de escolher um nome~$\metav{c}$ que não ocorra na prova (até então), para $\exists$I podemos — e em geral devemos — escolher um nome~$\metav{c}$ que já ocorra na prova. Assim como no caso de $\eor$I, muitas vezes não é claro qual $\metav{c}$ funcionará; portanto, para evitar retrocessos, você deve trabalhar para trás a partir de sentenças quantificadas existencialmente somente depois de aplicar todas as outras estratégias.

Em contraste, trabalhar \emph{para a frente} a partir de sentenças $\exists \metav{x}\, \atom{\metav{A}}{\metav{x}}$ em geral sempre funciona, e você não precisará voltar atrás. Trabalhar para a frente a partir de uma sentença quantificada existencialmente leva em conta não apenas $\exists \metav{x}\, \atom{\metav{A}}{\metav{x}}$, mas também qualquer sentença $\metav{B}$ que você queira provar. É preciso estabelecer uma subprova acima de $\metav{B}$, na qual $\metav{B}$ seja a última linha, e usar como hipótese uma instância de substituição $\atom{\metav{A}}{\metav{c}}$ de $\exists \metav{x}\, \atom{\metav{A}}{\metav{x}}$. Para garantir que a condição sobre $\metav{c}$ que rege $\exists$E seja satisfeita, escolha um nome~$\metav{c}$ que ainda não ocorra na prova.
\begin{fitchproof}
	\ellipsesline
	\have[m]{m}{\exists \metav{x}\, \atom{\metav{A}}{\metav{x}}}
	\ellipsesline
	\open
	\hypo[n]{n}{\atom{\metav{A}}{\metav{c}}}\AS
	\ellipsesline
	\have[k]{k}{\metav{B}}
	\close
	\have{e}{\metav{B}}\Ee{m,n-k}
\end{fitchproof}
Você então continuará com o objetivo de provar $\metav{B}$, mas agora dentro de uma subprova na qual tem uma sentença adicional com que trabalhar, a saber, $\atom{\metav{A}}{\metav{c}}$.

Por fim, trabalhar para a frente a partir de $\forall \metav{x}\, \atom{\metav{A}}{\metav{x}}$ significa que você sempre pode escrever $\atom{\metav{A}}{\metav{c}}$ e justificá-la usando $\forall$E, para qualquer nome~$\metav{c}$. É claro que você não deve fazer isso de qualquer maneira. Somente alguns nomes~$\metav{c}$ ajudarão na tarefa de provar a sentença-alvo em que está trabalhando. Portanto, assim como ao trabalhar para trás a partir de $\exists \metav{x}\, \atom{\metav{A}}{\metav{x}}$, você deve trabalhar para a frente a partir de $\forall \metav{x}\, \atom{\metav{A}}{\metav{x}}$ somente depois de aplicar todas as outras estratégias.

Consideremos como exemplo o argumento $\forall x(\atom{A}{x} \eif B) \therefore \exists x\,\atom{A}{x} \eif B$. Para começar a construir uma prova, escrevemos a premissa no alto e a conclusão embaixo.
\begin{fitchproof}
\hypo{1}{\forall x(\atom{A}{x} \eif B)}\PR
\ellipsesline
\have[n]{7}{\exists x\,\atom{A}{x} \eif B}
\end{fitchproof}
As estratégias para os conectivos da LVF continuam valendo, e você deve aplicá-las na mesma ordem: primeiro trabalhe para trás a partir de condicionais, sentenças negadas, conjunções e agora também quantificadores universais; depois trabalhe para a frente a partir de disjunções e agora também quantificadores existenciais; só então tente aplicar $\eif$E, $\enot$E, $\lor$I, $\forall$E ou $\exists$I. No nosso caso, isso significa trabalhar para trás a partir da conclusão:
\begin{fitchproof}
	\hypo{1}{\forall x(\atom{A}{x} \eif B)}\PR
	\open
	\hypo{2}{\exists x\,\atom{A}{x}}\AS
	\ellipsesline
	\have[n][-1]{6}{B}
	\close
	\have[n]{7}{\exists x\,\atom{A}{x} \eif B}\ci{2-(6)}
\end{fitchproof}
Nosso próximo passo deve ser trabalhar para a frente a partir de $\exists x\,\atom{A}{x}$ na linha~$2$. Para isso, precisamos escolher um nome que ainda não esteja na prova. Como não aparece nenhum nome, podemos escolher qualquer um, digamos~\textquotedblleft $d$\textquotedblright{}.
\begin{fitchproof}
	\hypo{1}{\forall x(\atom{A}{x} \eif B)}\PR
	\open
	\hypo{2}{\exists x\,\atom{A}{x}}\AS
	\open
	\hypo{3}{\atom{A}{d}}\AS
	\ellipsesline
	\have[n][-2]{5}{B}
	\close
	\have[n][-1]{6}{B}\Ee{2,3-(5)}
	\close
	\have[n]{7}{\exists x\,\atom{A}{x} \eif B}\ci{2-(6)}
\end{fitchproof}
Agora esgotamos nossas estratégias principais, e é hora de trabalhar para a frente a partir da premissa $\forall x(\atom{A}{x} \eif B)$. Aplicar $\forall$E significa que podemos justificar qualquer instância de $A(\metav{c}) \eif B$, independentemente de qual $\metav{c}$ escolhermos. É claro que será melhor escolher $d$, pois isso nos dará $\atom{A}{d} \eif B$. Então podemos aplicar $\eif$E para justificar~$B$, concluindo a prova.
\begin{fitchproof}
	\hypo{1}{\forall x(\atom{A}{x} \eif B)}\PR
	\open
	\hypo{2}{\exists x\,\atom{A}{x}}\AS
	\open
	\hypo{3}{\atom{A}{d}}\AS
\have{4}{\atom{A}{d} \eif B}\Ae{1}
	\have{5}{B}\ce{4,3}
	\close
	\have{6}{B}\Ee{2,3-5}
	\close
	\have{7}{\exists x\,\atom{A}{x} \eif B}\ci{2-6}
\end{fitchproof}

Agora vamos construir uma prova da recíproca. Começamos com
\begin{fitchproof}
	\hypo{1}{\exists x\,\atom{A}{x} \eif B}\PR
	\ellipsesline
	\have[n]{7}{\forall x(\atom{A}{x} \eif B)}
\end{fitchproof}
Observe que a premissa é um condicional, não uma sentença quantificada existencialmente; portanto, não devemos (ainda) trabalhar para a frente a partir dela. Trabalhar para trás a partir da conclusão, $\forall x(\atom{A}{x} \eif B)$, nos leva a procurar uma prova de $\atom{A}{d} \eif B$:
\begin{fitchproof}
	\hypo{1}{\exists x\,\atom{A}{x} \eif B}\PR
	\ellipsesline
	\have[n][-1]{6}{\atom{A}{d} \eif B}
	\have[n]{7}{\forall x(\atom{A}{x} \eif B)}\Ai{6}
\end{fitchproof}
E trabalhar para trás a partir de $\atom{A}{d} \eif B$ significa estabelecer uma subprova com $\atom{A}{d}$ como hipótese e $B$ como última linha:
\begin{fitchproof}
	\hypo{1}{\exists x\,\atom{A}{x} \eif B}\PR
	\open
	\hypo{2}{\atom{A}{d}}\AS
	\ellipsesline
	\have[n][-2]{5}{B}
	\close
	\have[n][-1]{6}{\atom{A}{d} \eif B}\ci{2-(5)}
	\have[n]{7}{\forall x(\atom{A}{x} \eif B)}\Ai{6}
\end{fitchproof}
Agora podemos trabalhar para a frente a partir da premissa na linha~$1$. Ela é um condicional, e seu consequente é justamente a sentença~$B$ que estamos tentando justificar. Portanto, devemos procurar uma prova de seu antecedente, $\exists x\,\atom{A}{x}$. É claro que isso já está prontamente disponível, por $\exists$I a partir da linha~$2$, e terminamos:
\begin{fitchproof}
	\hypo{1}{\exists x\,\atom{A}{x} \eif B}\PR
	\open
	\hypo{2}{\atom{A}{d}}\AS
	\have{3}{\exists x\,\atom{A}{x}}\Ei{2}
	\have{5}{B}\ce{1,3}
	\close
	\have{6}{\atom{A}{d} \eif B}\ci{2-5}
	\have{7}{\forall x(\atom{A}{x} \eif B)}\Ai{6}
\end{fitchproof}

\practiceproblems

\problempart
Use as estratégias para encontrar provas para cada um dos argumentos e teoremas a seguir:
\begin{compactlist}
\item $A \eif \forall x\,\atom{B}{x} \therefore \forall x(A \eif \atom{B}{x})$
\item $\exists x(A \eif \atom{B}{x}) \therefore A \eif \exists x\, \atom{B}{x}$
\item $\forall x(\atom{A}{x} \eand \atom{B}{x}) \eiff (\forall x\,\atom{A}{x} \eand \forall x\,\atom{B}{x})$
\item $\exists x(\atom{A}{x} \eor \atom{B}{x}) \eiff (\exists x\,\atom{A}{x} \eor \exists x\,\atom{B}{x})$
\item $A \eor \forall x\,\atom{B}{x}) \therefore \forall x(A \eor \atom{B}{x})$
\item $\forall x(\atom{A}{x} \eif B) \therefore \exists x\,\atom{A}{x} \eif B$
\item $\exists x(\atom{A}{x} \eif B) \therefore \forall x\,\atom{A}{x} \eif B$
\item $\forall x(\atom{A}{x} \eif \exists y\,\atom{A}{y})$
\end{compactlist}
Use apenas as regras básicas da LVF além das regras básicas para os quantificadores.

\problempart
Use as estratégias para encontrar provas para cada um dos argumentos e teoremas a seguir:
\begin{compactlist}
\item $\forall x\,\atom{R}{x,x} \therefore \forall x\,\exists y\,\atom{R}{x,y}$
\item $\forall x\,\forall y\,\forall z[(\atom{R}{x,y} \eand \atom{R}{y,z}) \eif \atom{R}{x,z}]$ \\
$\therefore \forall x\,\forall y[\atom{R}{x,y} \eif \forall z(\atom{R}{y,z} \eif \atom{R}{x,z})]$
\item $\forall x\,\forall y\,\forall z[(\atom{R}{x,y} \eand \atom{R}{y,z}) \eif \atom{R}{x,z}],$\\ $\forall x\,\forall y(\atom{R}{x,y} \eif \atom{R}{y, x})$ \\ $\therefore \forall x\,\forall y\,\forall z[(\atom{R}{x,y} \eand \atom{R}{x,z}) \eif \atom{R}{y,z}]$
\item $\forall x\,\forall y(\atom{R}{x,y} \eif \atom{R}{y, x})$ \\$\therefore \forall x\,\forall y\,\forall z[(\atom{R}{x,y} \eand \atom{R}{x,z}) \eif \exists u(\atom{R}{y,u} \eand \atom{R}{z,u})]$
\item $\enot \exists x\,\forall y (\atom{A}{x, y} \eiff \lnot\atom{A}{y, y})$
\end{compactlist}

\problempart
Use as estratégias para encontrar provas para cada um dos argumentos e teoremas a seguir:
\begin{compactlist}
\item $\forall x\,\atom{A}{x} \eif B \therefore \exists x(\atom{A}{x} \eif B)$
\item $A \eif \exists x\, \atom{B}{x} \therefore \exists x(A \eif \atom{B}{x})$
\item $\forall x(A \eor \atom{B}{x}) \therefore A \eor \forall x\,\atom{B}{x})$
\item $\exists x(\atom{A}{x} \eif \forall y\,\atom{A}{y})$
\item $\exists x(\exists y\,\atom{A}{y} \eif \atom{A}{x})$
\end{compactlist}
Estes exigem o uso de IP. Use apenas as regras básicas da LVF além das regras básicas para os quantificadores.

\chapter{Conversão de quantificadores}\label{s:CQ}

Nesta seção, acrescentaremos algumas regras às regras básicas da seção anterior. Elas regem a interação entre quantificadores e negação.

Em \cref{s:FOLBuildingBlocks}, observamos que $\enot\exists x\metav{A}$ é logicamente equivalente a $\forall x\,\enot\metav{A}$. Acrescentaremos ao nosso sistema de provas algumas regras que regem isso. Em particular, acrescentamos:
	\factoidbox{
	\begin{fitchproof}
		\have[m]{a}{\forall \metav{x}\, \enot\metav{A}}
		\have[\ ]{con}{\enot \exists \metav{x}\, \metav{A}}\cq{a}
	\end{fitchproof}}
e
\factoidbox{
	\begin{fitchproof}
		\have[m]{a}{ \enot \exists \metav{x}\, \metav{A}}
		\have[\ ]{con}{\forall \metav{x}\, \enot \metav{A}}\cq{a}
	\end{fitchproof}}
Do mesmo modo, acrescentamos:
\factoidbox{
	\begin{fitchproof}
		\have[m]{a}{\exists \metav{x}\, \enot \metav{A}}
		\have[\ ]{con}{\enot \forall \metav{x}\, \metav{A}}\cq{a}
	\end{fitchproof}}
e
\factoidbox{
	\begin{fitchproof}
		\have[m]{a}{\enot \forall \metav{x}\, \metav{A}}
		\have[\ ]{con}{\exists \metav{x}\, \enot \metav{A}}\cq{a}
	\end{fitchproof}}

\practiceproblems
\problempart
Mostre, em cada caso, que as sentenças são inconsistentes:
\begin{compactlist}
\item $\atom{S}{a}\eif \atom{T}{m}, \atom{T}{m} \eif \atom{S}{a}, \atom{T}{m} \eand \enot \atom{S}{a}$
\item $\enot\exists x\,\atom{R}{x,a}, \forall x\, \forall y\,\atom{R}{y,x}$
\item $\enot\exists x\, \exists y\,\atom{L}{x,y}, \atom{L}{a,a}$
\item $\forall x(\atom{P}{x} \eif \atom{Q}{x}), \forall z(\atom{P}{z} \eif \atom{R}{z}), \forall y\,\atom{P}{y}, \enot \atom{Q}{a} \eand \enot \atom{R}{b}$
\end{compactlist}

\problempart
Mostre que cada par de sentenças é interderivável:
\begin{compactlist}
\item $\forall x (\atom{A}{x}\eif \enot \atom{B}{x}), \enot\exists x(\atom{A}{x} \eand \atom{B}{x})$
\item $\forall x (\enot\atom{A}{x}\eif \atom{B}{d}), \forall x\,\atom{A}{x} \eor \atom{B}{d}$
\end{compactlist}

\problempart
Em \cref{s:MoreMonadic}, consideramos o que acontece quando movemos quantificadores \textquotedblleft através\textquotedblright{} de vários operadores lógicos. Mostre que cada par de sentenças é interderivável:
\begin{compactlist}
\item $\forall x(\atom{F}{x} \eand \atom{G}{a}), \forall x\,\atom{F}{x} \eand \atom{G}{a}$
\item $\exists x(\atom{F}{x} \eor \atom{G}{a}), \exists x\,\atom{F}{x} \eor \atom{G}{a}$
\item $\forall x(\atom{G}{a} \eif \atom{F}{x}), \atom{G}{a} \eif \forall x\,\atom{F}{x}$
\item $\forall x(\atom{F}{x} \eif \atom{G}{a}), \exists x\,\atom{F}{x} \eif \atom{G}{a}$
\item $\exists x(\atom{G}{a} \eif \atom{F}{x}), \atom{G}{a} \eif \exists x\,\atom{F}{x}$
\item $\exists x(\atom{F}{x} \eif \atom{G}{a}), \forall x\,\atom{F}{x} \eif \atom{G}{a}$
\end{compactlist}
OBS.: a variável \textquotedblleft $x$\textquotedblright{} não ocorre em \textquotedblleft $\atom{G}{a}$\textquotedblright{}. Quando todos os quantificadores ocorrem no início de uma sentença, dizemos que ela está na \emph{forma normal prenex}. Essas equivalências às vezes são chamadas de \emph{regras de prenexação}, pois fornecem um meio de colocar qualquer sentença na forma normal prenex.


\chapter{Regras para a identidade}
Em \cref{s:Interpretations}, mencionamos a tese filosoficamente controversa da \emph{identidade dos indiscerníveis}. Essa é a afirmação de que objetos indiscerníveis em todos os aspectos são, de fato, idênticos uns aos outros. Também mencionamos que não aderiremos a essa tese. Segue-se que, por mais que você aprenda sobre dois objetos, não podemos provar que são idênticos. A menos, é claro, que você descubra que os dois objetos são de fato idênticos; mas, nesse caso, a prova dificilmente será muito esclarecedora.

O ponto geral, porém, é que \emph{nenhuma sentença} que não contenha de antemão o predicado de identidade poderia justificar uma inferência para \textquotedblleft $a=b$\textquotedblright{}. Portanto, nossa regra de introdução da identidade não pode permitir inferir uma afirmação de identidade contendo dois nomes \emph{diferentes}.

No entanto, todo objeto é idêntico a si mesmo. Portanto, nenhuma premissa é necessária para concluir que algo é idêntico a si mesmo. Assim, esta será a regra de introdução da identidade:
\factoidbox{
\begin{fitchproof}
	\have[\ \,\,\,]{x}{\metav{c}=\metav{c}} \by{=I}{}
\end{fitchproof}}
Observe que essa regra não exige referência a nenhuma linha anterior da prova. Para qualquer nome $\metav{c}$, você pode escrever $\metav{c}=\metav{c}$ em qualquer linha, tendo apenas a regra {=}I como justificativa.

Nossa regra de eliminação é mais divertida. Se você estabeleceu \textquotedblleft $a=b$\textquotedblright{}, então tudo o que é verdadeiro do objeto nomeado por \textquotedblleft $a$\textquotedblright{} também deve ser verdadeiro do objeto nomeado por \textquotedblleft $b$\textquotedblright{}. Para qualquer sentença que contenha \textquotedblleft $a$\textquotedblright{}, você pode substituir algumas ou todas as ocorrências de \textquotedblleft $a$\textquotedblright{} por \textquotedblleft $b$\textquotedblright{} e produzir uma sentença equivalente. Por exemplo, de \textquotedblleft $\atom{R}{a,a}$\textquotedblright{} e \textquotedblleft $a=b$\textquotedblright{}, você está justificado a inferir \textquotedblleft $\atom{R}{a,b}$\textquotedblright{}, \textquotedblleft $\atom{R}{b,a}$\textquotedblright{} ou \textquotedblleft $\atom{R}{b,b}$\textquotedblright{}. Mais geralmente:
\factoidbox{\begin{fitchproof}
	\have[m]{e}{\metav{a}=\metav{b}}
	\have[n]{a}{\metav{A}(\ldots \metav{a} \ldots \metav{a}\ldots)}
	\have[\ ]{ea1}{\metav{A}(\ldots \metav{b} \ldots \metav{a}\ldots)} \by{=E}{e,a}
\end{fitchproof}}
A notação aqui é a mesma da usada para $\exists$I. Assim, $\metav{A}(\ldots \metav{a} \ldots \metav{a}\ldots)$ é uma fórmula que contém o nome $\metav{a}$, e $\metav{A}(\ldots \metav{b} \ldots \metav{a}\ldots)$ é uma fórmula obtida substituindo uma ou mais ocorrências do nome $\metav{a}$ pelo nome $\metav{b}$. As linhas $m$ e $n$ podem ocorrer em qualquer ordem e não precisam ser adjacentes, mas sempre citamos primeiro a afirmação de identidade. Simetricamente, permitimos:
\factoidbox{\begin{fitchproof}
	\have[m]{e}{\metav{a}=\metav{b}}
	\have[n]{a}{\metav{A}(\ldots \metav{b} \ldots \metav{b}\ldots)}
	\have[\ ]{ea2}{\metav{A}(\ldots \metav{a} \ldots \metav{b}\ldots)} \by{=E}{e,a}
\end{fitchproof}}
Essa regra às vezes é chamada de \emph{lei de Leibniz}, em homenagem a Gottfried Leibniz.

Para ver as regras em ação, provaremos alguns resultados rápidos. Primeiro, provaremos que a identidade é \emph{simétrica}:
\begin{fitchproof}
	\open
		\hypo{ab}{a = b}\AS
		\have{aa}{a = a}\by{=I}{}
		\have{ba}{b = a}\by{=E}{ab, aa}
	\close
	\have{abba}{a = b \eif b =a}\ci{ab-ba}
	\have{ayya}{\forall y (a = y \eif y = a)}\Ai{abba}
	\have{xyyx}{\forall x\, \forall y (x = y \eif y = x)}\Ai{ayya}
\end{fitchproof}
Obtemos a linha~$3$ substituindo uma ocorrência de \textquotedblleft $a$\textquotedblright{} na linha~$2$ por uma ocorrência de \textquotedblleft $b$\textquotedblright{}; isso é justificado por \textquotedblleft $a=b$\textquotedblright{}.

Segundo, provaremos que a identidade é \emph{transitiva}:
\begin{fitchproof}
	\open
		\hypo{abc}{a = b \eand b = c}\AS
		\have{ab}{a = b}\ae{abc}
		\have{bc}{b = c}\ae{abc}
		\have{ac}{a = c}\by{=E}{ab, bc}
	\close
	\have{con}{(a = b \eand b =c) \eif a = c}\ci{abc-ac}
	\have{conz}{\forall z((a = b \eand b = z) \eif a = z)}\Ai{con}
	\have{cony}{\forall y\,\forall z((a = y \eand y = z) \eif a = z)}\Ai{conz}
	\have{conx}{\forall x\,\forall y \forall z((x = y \eand y = z) \eif x = z)}\Ai{cony}
\end{fitchproof}
Obtemos a linha~$4$ substituindo \textquotedblleft $b$\textquotedblright{} na linha~$3$ por \textquotedblleft $a$\textquotedblright{}; isso é justificado por \textquotedblleft $a=b$\textquotedblright{}.

\practiceproblems
\problempart
\label{pr.identity}
Para cada uma das afirmações a seguir, forneça uma prova na LPO que mostre
que ela é verdadeira.
\begin{compactlist}
\item $\atom{P}{a} \eor \atom{Q}{b}, \atom{Q}{b} \eif b=c, \enot\atom{P}{a} \proves \atom{Q}{c}$
\item $m=n \eor n=o, \atom{A}{n} \proves \atom{A}{m} \eor \atom{A}{o}$
\item $\forall x\ x=m, \atom{R}{m,a} \proves \exists x\,\atom{R}{x,x}$
\item $\forall x\,\forall y(\atom{R}{x,y} \eif x=y)\proves \atom{R}{a,b} \eif \atom{R}{b,a}$
\item $\enot \exists x\enot\, {x = m} \proves \forall x\,\forall y (\atom{P}{x} \eif \atom{P}{y})$
\item $\exists x\,\atom{J}{x}, \exists x\, \enot\atom{J}{x}\proves \exists x\, \exists y\, \enot\, {x = y}$
\item $\forall x(x=n \eiff \atom{M}{x}), \forall x(\atom{O}{x} \eor \enot \atom{M}{x})\proves \atom{O}{n}$
\item $\exists x\,\atom{D}{x}, \forall x(x=p \eiff \atom{D}{x})\proves \atom{D}{p}$
\item $\exists x\bigl[(\atom{K}{x} \eand \forall y(\atom{K}{y} \eif x=y)) \eand \atom{B}{x}\bigr], \atom{K}{d} \proves \atom{B}{d}$
\item $\proves \atom{P}{a} \eif \forall x(\atom{P}{x} \eor \enot\, {x = a})$
\end{compactlist}

\problempart
Mostre que as seguintes sentenças são interderiváveis:
\begin{itemize}
\item $\exists x \bigl([\atom{F}{x} \eand \forall y (\atom{F}{y} \eif x = y)] \eand x = n\bigr)$
\item $\atom{F}{n} \eand \forall y (\atom{F}{y} \eif n= y)$
\end{itemize}
E, portanto, ambas têm uma boa pretensão de simbolizar a sentença em português \textquotedblleft Nick é o~$F$\textquotedblright{}.

\problempart
Em \cref{sec.identity}, afirmamos que as seguintes são simbolizações logicamente equivalentes da sentença em português \textquotedblleft existe exatamente um $F$\textquotedblright{}:
\begin{itemize}
\item $\exists x\,\atom{F}{x} \eand \forall x\, \forall y \bigl[(\atom{F}{x} \eand \atom{F}{y}) \eif x = y\bigr]$
\item $\exists x \bigl[\atom{F}{x} \eand \forall y (\atom{F}{y} \eif x = y)\bigr]$
\item $\exists x\, \forall y (\atom{F}{y} \eiff x = y)$
\end{itemize}
Mostre que todas são interderiváveis. (\emph{Dica}: para mostrar que três afirmações são interderiváveis, basta mostrar que a primeira prova a segunda, a segunda prova a terceira e a terceira prova a primeira; pense no motivo.)


\
\problempart
Simbolize o seguinte argumento
	\begin{earg}
		\item Existe exatamente um $F$.
		\item Existe exatamente um $G$.
		\item Nada é simultaneamente $F$ e~$G$.
		\item[\texttherefore] Há exatamente duas coisas que são $F$ ou~$G$.
	\end{earg}
E ofereça uma prova dele.
%\begin{itemize}
%\item  $\exists x \bigl[\atom{F}{x} \eand \forall y (\atom{F}{y} \eif x = y)\bigr], \exists x \bigl[\atom{G}{x} \eand \forall y (\atom{G}{y} \eif x = y)\bigr], \forall x (\enot\atom{F}{x} \eor \enot \atom{G}{x}) \proves \exists x\, \exists y \bigl[\enot\, {x = y} \eand \forall z ((\atom{F}{z} \eor \atom{G}{z}) \eif (x = y \eor x = z))\bigr]$
%\end{itemize}




\chapter{Regras derivadas}\label{s:DerivedFOL}
Assim como no caso da LVF, primeiro introduzimos algumas regras para a LPO como básicas (em \cref{s:BasicFOL}) e depois acrescentamos outras regras para a conversão de quantificadores (em \cref{s:CQ}). Na verdade, as regras CQ devem ser consideradas regras \emph{derivadas}, pois podem ser derivadas das regras \emph{básicas} de \cref{s:BasicFOL}. (O ponto aqui é o mesmo de \cref{s:Derived}.) Eis uma justificativa para a primeira regra CQ:
\begin{fitchproof}
	\have[m]{An}{\forall x\, \enot \atom{A}{x}}
	\open
		\hypo[m+1]{E}{\exists x\, \atom{A}{x}}\AS
		\open
			\hypo[m+2]{c}{\atom{A}{c}}\AS
			\have[m+3]{nc}{\enot \atom{A}{c}}\Ae{An}
			\have[m+4]{red}{\ered}\ne{nc,c}
		\close
		\have[m+5]{red2}{\ered}\Ee{E,(c)-(red)}
	\close
	\have[m+6]{dada}{\enot \exists x\, \atom{A}{x}}\ni{(E)-(red2)}
\end{fitchproof}
%You will note that on line 3 I have written `for $\exists$E'. This is not technically a part of the proof. It is just a reminder---to me and to you---of why I have bothered to introduce `$\enot \atom{A}{c}$' out of the blue. You might find it helpful to add similar annotations to assumptions when performing proofs. But do not add annotations on lines other than assumptions: the proof requires its own citation, and your annotations will clutter it.
Eis uma justificativa para a terceira regra CQ:
\begin{fitchproof}
	\have[m]{nEna}{\exists x\, \enot \atom{A}{x}}
	\open
		\hypo[m+1]{Aa}{\forall x\, \atom{A}{x}}\AS
		\open
			\hypo[m+2]{nac}{\enot \atom{A}{c}}\AS
			\have[m+3]{a}{\atom{A}{c}}\Ae{Aa}
			\have[m+4]{con}{\ered}\ne{nac,a}
		\close
		\have[m+5]{con1}{\ered}\Ee{nEna, (nac)-(con)}
	\close
	\have[m+6]{dada}{\enot \forall x\, \atom{A}{x}}\ni{(Aa)-(con1)}
\end{fitchproof}
Isso explica por que as regras CQ podem ser tratadas como derivadas. Justificativas semelhantes podem ser oferecidas para as outras duas regras CQ.

\practiceproblems

\problempart
Ofereça provas que justifiquem a inclusão da segunda e da quarta regras CQ como regras derivadas.



\chapter{Provas e semântica}
Usamos dois símbolos de consequência diferentes neste livro. Este:
$$\metav{A}_1, \metav{A}_2, \ldots, \metav{A}_n \proves \metav{C}$$
significa que há alguma prova que termina com $\metav{C}$ e cujas únicas hipóteses não descarregadas estão entre $\metav{A}_1, \metav{A}_2, \ldots, \metav{A}_n$. Esta é uma \emph{noção teórica da prova}. Em contraste, este:
$$\metav{A}_1, \metav{A}_2, \ldots, \metav{A}_n \entails \metav{C}$$
significa que nenhuma valoração (ou interpretação) torna verdadeiras todas as sentenças $\metav{A}_1, \metav{A}_2, \ldots, \metav{A}_n$ e~$\metav{C}$ falsa. Isso diz respeito às atribuições de verdade e falsidade às sentenças. É uma \emph{noção semântica}.

Não se pode enfatizar o suficiente que essas são noções diferentes. Mas podemos enfatizar um pouco mais: \emph{elas são noções diferentes}.

Depois de internalizar este ponto, continue a leitura.

Embora nossas noções semântica e teórica da prova sejam diferentes, há uma conexão profunda entre elas. Para explicar essa conexão, começaremos considerando a relação entre validades e teoremas.

Para mostrar que uma sentença é um teorema, basta produzir uma prova. É verdade que pode ser difícil produzir uma prova de vinte linhas, mas não é tão difícil verificar cada linha e confirmar que é legítima; e, se cada linha da prova for individualmente legítima, então a prova inteira será legítima. Mostrar que uma sentença é uma validade, porém, exige raciocinar sobre todas as interpretações possíveis. Dada a escolha entre mostrar que uma sentença é um teorema e mostrar que é uma validade, será mais fácil mostrar que é um teorema.

Inversamente, mostrar que uma sentença \emph{não} é um teorema é difícil. Teríamos de raciocinar sobre todas as provas (possíveis). Isso é muito difícil. No entanto, para mostrar que uma sentença não é uma validade, basta construir uma interpretação na qual a sentença seja falsa. É verdade que pode ser difícil encontrar a interpretação; mas, uma vez encontrada, é relativamente simples verificar que valor de verdade ela atribui à sentença. Dada a escolha entre mostrar que uma sentença não é um teorema e mostrar que não é uma validade, seria mais fácil mostrar que não é uma validade.

Felizmente, \emph{uma sentença é um teorema se, e somente se, é uma validade}. Como resultado, se fornecermos uma prova de $\metav{A}$ sem hipóteses e, assim, mostrarmos que $\metav{A}$ é um teorema, isto é, ${}\proves \metav{A}$, podemos inferir legitimamente que $\metav{A}$ é uma validade, isto é, $\entails\metav{A}$. Da mesma forma, se construirmos uma interpretação na qual $\metav{A}$ é falsa e, assim, mostrarmos que não é uma validade, isto é, $\nentails \metav{A}$, segue-se que $\metav{A}$ não é um teorema, isto é, $\nproves \metav{A}$.

De modo mais geral, temos o seguinte resultado poderoso:
\factoidbox{$\metav{A}_1, \metav{A}_2, \ldots, \metav{A}_n
\proves\metav{B}$ \textbf{\ifeff} $\metav{A}_1, \metav{A}_2, \ldots,
\metav{A}_n \entails\metav{B}$}
Isso mostra que, embora demonstrabilidade e consequência lógica sejam noções \emph{diferentes}, elas são extensionalmente equivalentes. Assim:
	\factoidbox{\begin{factoiditems}
		\item Um argumento é \emph{válido} \ifeff{} \emph{a conclusão pode ser provada a partir das premissas}.
		\item Uma sentença é uma \emph{validade} \ifeff{} é um \emph{teorema}.
		\item Duas sentenças são \emph{equivalentes} \ifeff{} são
		\emph{interderiváveis}.
		\item Sentenças são \emph{satisfatíveis em conjunto} \ifeff{} são
		\emph{consistentes em conjunto}.
	\end{factoiditems}}
Por esse motivo, você pode escolher quando pensar em termos de provas e quando pensar em termos de valorações/interpretações, fazendo o que for mais fácil para uma dada tarefa. A tabela da página seguinte resume o que costuma ser mais fácil.

É intuitivo que demonstrabilidade e consequência semântica devam coincidir. Mas — repitamos isto — não se deixe enganar pela semelhança dos símbolos \textquotedblleft $\entails$\textquotedblright{} e \textquotedblleft $\proves$\textquotedblright{}. Esses dois símbolos têm significados muito diferentes. O fato de demonstrabilidade e consequência semântica coincidirem não é um resultado fácil de obter.

De fato, demonstrar que demonstrabilidade e consequência semântica coincidem é, decisivamente, o ponto em que a lógica introdutória se torna lógica intermediária.

\ifHTMLtarget
\begin{table}
\else
\begin{sidewaystable}\small
\fi
	\begin{center}
\begin{tabular*}{\textwidth}{p{.25\textheight}p{.325\textheight}p{.325\textheight}}
 & \textbf{Sim}  & \textbf{Não}\\[12pt]
A sentença $\metav{A}$ é uma \textbf{validade}?
& dê uma prova que mostre $\proves\metav{A}$
& dê uma interpretação na qual $\metav{A}$ seja falsa\\[12pt]
A sentença $\metav{A}$ é uma \textbf{contradição}?
& dê uma prova que mostre $\metav{A} \proves \ered$
& dê uma interpretação na qual $\metav{A}$ seja verdadeira\\[12pt]
%Is \metav{A} contingent? &
%give two interpretations, one in which \metav{A} is true and another in which \metav{A} is false & give a proof which either shows $\proves\metav{A}$ or $\proves\enot\metav{A}$\\
%\\
As sentenças $\metav{A}$ e $\metav{B}$ são \textbf{equivalentes}?
& dê duas provas, uma de $\metav{A}\proves\metav{B}$ e outra de $\metav{B}\proves\metav{A}$
& dê uma interpretação na qual $\metav{A}$ e $\metav{B}$ tenham valores de verdade diferentes\\[12pt]
As sentenças $\metav{A}_1, \metav{A}_2, \ldots, \metav{A}_n$ são \textbf{conjuntamente satisfatíveis}?
& dê uma interpretação na qual todas as sentenças $\metav{A}_1, \metav{A}_2, \ldots, \metav{A}_n$ sejam verdadeiras
& prove uma contradição a partir das hipóteses $\metav{A}_1, \metav{A}_2, \ldots, \metav{A}_n$\\[12pt]
O argumento $\metav{A}_1, \metav{A}_2, \ldots, \metav{A}_n \therefore \metav{C}$ é \textbf{válido}?
& dê uma prova com as hipóteses $\metav{A}_1, \metav{A}_2, \ldots, \metav{A}_n$ e conclusão $\metav{C}$
& dê uma interpretação na qual cada uma de $\metav{A}_1, \metav{A}_2, \ldots, \metav{A}_n$ seja verdadeira e $\metav{C}$ seja falsa\\
\end{tabular*}
\end{center}
\ifHTMLtarget
\end{table}
\else
\end{sidewaystable}
\fi

